% Generated from project outputs. Do not edit by hand.
\begin{table}[!htbp]
\centering
\begin{threeparttable}
\caption{Including in-kind by subsample: 2016 moments}
\label{tab:annual-2016-y2-moments}
\small
\setlength{\tabcolsep}{4pt}
\renewcommand{\arraystretch}{1.15}
\begin{tabularx}{\linewidth}{@{}Xrrrrr@{}}
\toprule
Sample & N & Mean & SD & Min & Max \\
\midrule
General & 94 & 2,761.1 & 1,666.5 & 244.0 & 7,051.6 \\
Men & 94 & 3,064.7 & 1,901.6 & 365.4 & 9,340.4 \\
Women & 94 & 2,205.8 & 1,408.3 & 85.4 & 5,690.6 \\
Urban & 94 & 3,021.0 & 1,737.1 & 369.6 & 7,703.2 \\
Rural & 94 & 2,085.7 & 1,339.7 & 96.2 & 8,577.9 \\
Indigenous & 92 & 2,020.1 & 1,295.6 & 297.2 & 5,977.6 \\
Non-Indigenous & 94 & 2,971.6 & 1,676.6 & 169.4 & 7,285.3 \\
Bottom 99 percent & 94 & 2,503.4 & 1,316.2 & 244.0 & 5,399.1 \\
Bottom 90 percent & 94 & 2,078.0 & 943.1 & 149.7 & 4,935.3 \\
Bottom 50 percent & 94 & 1,143.8 & 424.9 & 135.9 & 2,685.7 \\
\bottomrule
\end{tabularx}
\begin{tablenotes}[flushleft]
\footnotesize
\item Monthly quetzales. Unweighted distribution of survey-weighted cohort means before transition matching. N counts cohort cells, not individuals. SD is dispersion, not a standard error. Subsamples overlap. Cumulative income definitions and treatment of missing components follow the merging scripts.
\end{tablenotes}
\end{threeparttable}
\end{table}
